
\documentclass[12pt,a4paper]{article}

% Поддержка русского языка
\usepackage[T2A]{fontenc}
\usepackage[utf8]{inputenc}
\usepackage[russian]{babel}

% Математика
\usepackage{amsmath,amssymb,amsthm}

% Настройка страницы
\usepackage[a4paper,margin=2.5cm]{geometry}

\begin{document}
	
	\tableofcontents
	
	\newpage
	\section{Получебышёвский, растяжка вправо}
	
	Рассмотрим систему
	\begin{equation}
		\beta
		\begin{pmatrix}
			0\\
			1\\
			2x\\
			3x^2
		\end{pmatrix}
		=
		-\omega_1
		\begin{pmatrix}
			1\\
			-1\\
			1\\
			-1
		\end{pmatrix}
		+
		\omega_2
		\begin{pmatrix}
			1\\
			t_1\\
			t_1^2\\
			t_1^3
		\end{pmatrix}
		-
		(1-\omega_1-\omega_2)
		\begin{pmatrix}
			1\\
			t_2\\
			t_2^2\\
			t_2^3
		\end{pmatrix},
	\end{equation}
	
	а дополнительное условие теперь одно:
	\begin{equation}
		\frac{t_2-t_1}{t_1+1}=2.
		\tag{1}
	\end{equation}
	
	Необходимо найти
	\[
	(\omega_1,\omega_2,t_1,t_2,\beta).
	\]
	
	Будем считать $x$ заданным параметром.
	
	\subsection*{1. Дополнительное уравнение для $t_1,t_2$}
	
	Из
	\[
	\frac{t_2-t_1}{t_1+1}=2
	\]
	получаем
	\[
	t_2-t_1=2(t_1+1),
	\]
	откуда
	\[
	\boxed{t_2=3t_1+2}.
	\tag{2}
	\]
	
	Также необходимо
	\[
	\boxed{t_1\neq -1}.
	\tag{3}
	\]
	
	\subsection*{2. Разбиение векторного уравнения на скалярные}
	
	Получаем четыре уравнения:
	\begin{align}
		0&=-\omega_1+\omega_2-(1-\omega_1-\omega_2),
		\tag{4}\\
		\beta&=\omega_1+\omega_2t_1-(1-\omega_1-\omega_2)t_2,
		\tag{5}\\
		2\beta x&=-\omega_1+\omega_2t_1^2-(1-\omega_1-\omega_2)t_2^2,
		\tag{6}\\
		3\beta x^2&=\omega_1+\omega_2t_1^3-(1-\omega_1-\omega_2)t_2^3.
		\tag{7}
	\end{align}
	
	Именно последнее уравнение теперь даст дополнительное условие на $t_1$.
	
	\subsection*{3. Нахождение $\omega_2$}
	
	Из (4):
	\[
	0=-\omega_1+\omega_2-1+\omega_1+\omega_2,
	\]
	поэтому
	\[
	2\omega_2-1=0,
	\]
	то есть
	\[
	\boxed{\omega_2=\frac12}.
	\tag{8}
	\]
	
	Следовательно,
	\[
	1-\omega_1-\omega_2=\frac12-\omega_1.
	\tag{9}
	\]
	
	\subsection*{4. Нахождение $\beta$}
	
	Подставляем (8), (9) и $t_2=3t_1+2$ в (5):
	\[
	\beta
	=
	\omega_1+\frac{t_1}{2}
	-\left(\frac12-\omega_1\right)(3t_1+2).
	\]
	
	Раскрываем скобки:
	\[
	\beta
	=
	\omega_1+\frac{t_1}{2}
	-\frac{3t_1+2}{2}
	+\omega_1(3t_1+2).
	\]
	
	Получаем
	\[
	\beta=3\omega_1t_1+3\omega_1-t_1-1.
	\]
	
	Следовательно,
	\[
	\boxed{\beta=(t_1+1)(3\omega_1-1)}.
	\tag{10}
	\]
	
	\subsection*{5. Используем третье скалярное уравнение}
	
	Рассмотрим (6):
	\[
	2\beta x
	=
	-\omega_1+\frac{t_1^2}{2}
	-\left(\frac12-\omega_1\right)t_2^2.
	\]
	
	Так как
	\[
	t_2=3t_1+2,
	\]
	то
	\[
	t_2^2=9t_1^2+12t_1+4.
	\]
	
	Поэтому
	\[
	2\beta x
	=
	-\omega_1+\frac{t_1^2}{2}
	-\left(\frac12-\omega_1\right)
	(9t_1^2+12t_1+4).
	\]
	
	После раскрытия скобок:
	\[
	2\beta x
	=
	3\omega_1(3t_1^2+4t_1+1)
	-4t_1^2-6t_1-2.
	\]
	
	Замечаем:
	\[
	3t_1^2+4t_1+1=(t_1+1)(3t_1+1),
	\]
	а
	\[
	4t_1^2+6t_1+2
	=
	2(t_1+1)(2t_1+1).
	\]
	
	Следовательно,
	\[
	2\beta x
	=
	(t_1+1)
	\left[
	3\omega_1(3t_1+1)-2(2t_1+1)
	\right].
	\tag{11}
	\]
	
	С другой стороны, из (10):
	\[
	2\beta x
	=
	2x(t_1+1)(3\omega_1-1).
	\tag{12}
	\]
	
	Так как $t_1\neq -1$, сокращаем $t_1+1$:
	\[
	2x(3\omega_1-1)
	=
	3\omega_1(3t_1+1)-2(2t_1+1).
	\]
	
	Отсюда
	\[
	6x\omega_1-2x
	=
	9t_1\omega_1+3\omega_1-4t_1-2.
	\]
	
	Следовательно,
	\[
	\omega_1(6x-9t_1-3)
	=
	2x-4t_1-2,
	\]
	и
	\[
	\boxed{
		\omega_1=
		\frac{2(2t_1+1-x)}
		{3(3t_1+1-2x)}
	}.
	\tag{13}
	\]
	
	\subsection*{6. Теперь используем четвёртое уравнение}
	
	Вот здесь появляется отличие от предыдущего решения.
	
	Берём (7):
	\[
	3\beta x^2
	=
	\omega_1+\frac{t_1^3}{2}
	-\left(\frac12-\omega_1\right)t_2^3.
	\]
	
	Поскольку
	\[
	t_2=3t_1+2,
	\]
	имеем
	\[
	t_2^3
	=
	(3t_1+2)^3
	=
	27t_1^3+54t_1^2+36t_1+8.
	\]
	
	Поэтому
	\[
	3\beta x^2
	=
	\omega_1+\frac{t_1^3}{2}
	-\left(\frac12-\omega_1\right)
	(27t_1^3+54t_1^2+36t_1+8).
	\]
	
	Подставляем сюда выражения
	\[
	\beta=(t_1+1)(3\omega_1-1)
	\]
	и
	\[
	\omega_1=
	\frac{2(2t_1+1-x)}
	{3(3t_1+1-2x)}.
	\]
	
	После раскрытия скобок и приведения подобных членов получается
	\[
	\boxed{
		(t_1+1)^2
		\left(
		3t_1^2-(8x+2)t_1+3x^2-2x-2
		\right)=0.
	}
	\tag{14}
	\]
	
	Так как $t_1\neq -1$, имеем
	\[
	3t_1^2-(8x+2)t_1+3x^2-2x-2=0.
	\tag{15}
	\]
	
	\subsection*{7. Нахождение $t_1$}
	
	Дискриминант:
	\[
	D=(8x+2)^2-12(3x^2-2x-2).
	\]
	
	Раскрываем:
	\[
	D
	=
	64x^2+32x+4
	-36x^2+24x+24.
	\]
	
	Следовательно,
	\[
	D=28x^2+56x+28,
	\]
	то есть
	\[
	\boxed{D=28(x+1)^2}.
	\tag{16}
	\]
	
	Поэтому
	\[
	t_1
	=
	\frac{8x+2\pm\sqrt{28(x+1)^2}}{6}.
	\]
	
	Введём
	\[
	s\in\{-1,1\}.
	\]
	
	Тогда
	\[
	\boxed{
		t_1=
		\frac{1+4x+s\sqrt{7}\,(x+1)}{3}
	}.
	\tag{17}
	\]
	
	\subsection*{8. Нахождение $t_2$}
	
	Из (2):
	\[
	t_2=3t_1+2.
	\]
	
	Подставляем (17):
	\[
	t_2
	=
	1+4x+s\sqrt{7}(x+1)+2.
	\]
	
	Следовательно,
	\[
	\boxed{
		t_2=3+4x+s\sqrt{7}\,(x+1)
	}.
	\tag{18}
	\]
	
	\subsection*{9. Нахождение $\omega_1$}
	
	Используем (13):
	\[
	\omega_1=
	\frac{2(2t_1+1-x)}
	{3(3t_1+1-2x)}.
	\]
	
	Подставляем
	\[
	t_1=
	\frac{1+4x+s\sqrt{7}(x+1)}{3}.
	\]
	
	После упрощения:
	\[
	\boxed{
		\omega_1=
		\frac{2(4+s\sqrt{7})}{27}
	}.
	\tag{19}
	\]
	
	То есть
	\[
	\boxed{
		\omega_1=\frac{8+2s\sqrt{7}}{27}
	}.
	\]
	
	А
	\[
	\boxed{\omega_2=\frac12}.
	\tag{20}
	\]
	
	\subsection*{10. Нахождение $\beta$}
	
	Используем (10):
	\[
	\beta=(t_1+1)(3\omega_1-1).
	\]
	
	Для найденных значений получаем
	\[
	\boxed{
		\beta=
		\frac{(10+7s\sqrt{7})(x+1)}{27}
	}.
	\tag{21}
	\]
	
	\subsection*{11. Окончательный ответ}
	
	Таким образом, при
	\[
	s=\pm1
	\]
	получаем две ветви решения:
	\[
	\boxed{
		\omega_1=
		\frac{2(4+s\sqrt{7})}{27}
	},
	\]
	\[
	\boxed{
		\omega_2=\frac12
	},
	\]
	\[
	\boxed{
		t_1=
		\frac{1+4x+s\sqrt{7}(x+1)}{3}
	},
	\]
	\[
	\boxed{
		t_2=
		3+4x+s\sqrt{7}(x+1)
	},
	\]
	\[
	\boxed{
		\beta=
		\frac{(10+7s\sqrt{7})(x+1)}{27}
	}.
	\]
	
	Или компактно:
	\[
	\boxed{
		(\omega_1,\omega_2,t_1,t_2,\beta)
		=
		\left(
		\frac{2(4+s\sqrt{7})}{27},
		\frac12,
		\frac{1+4x+s\sqrt{7}(x+1)}{3},
		3+4x+s\sqrt{7}(x+1),
		\frac{(10+7s\sqrt{7})(x+1)}{27}
		\right),
		\qquad s=\pm1.
	}
	\]
	
	При этом $x=-1$ фактически не даёт допустимого решения: тогда обе ветви дают
	\[
	t_1=-1,
	\]
	что запрещено исходным знаменателем $t_1+1$.
	
	Таким образом, структура решения осталась почти такой же, как в исходном файле: четвёртое уравнение с кубами снова определяет $t_1$, только теперь из-за нового соотношения между $t_1$ и $t_2$ получается другое квадратное уравнение:
	\[
	\boxed{
		3t_1^2-(8x+2)t_1+3x^2-2x-2=0
	}.
	\]
	
	
	
	\newpage
	\section{Получебышёвский, растяжка влево}
	

		
		Рассмотрим систему
		\[
		\beta
		\begin{pmatrix}
			0\\
			1\\
			2x\\
			3x^2
		\end{pmatrix}
		=
		\omega_1
		\begin{pmatrix}
			1\\
			t_1\\
			t_1^2\\
			t_1^3
		\end{pmatrix}
		-
		\omega_2
		\begin{pmatrix}
			1\\
			t_2\\
			t_2^2\\
			t_2^3
		\end{pmatrix}
		+
		(1-\omega_1-\omega_2)
		\begin{pmatrix}
			1\\
			1\\
			1\\
			1
		\end{pmatrix},
		\]
		
		а дополнительное условие имеет вид
		\[
		\frac{1-t_2}{t_2-t_1}
		=
		\frac{1-\frac12}{\frac12+\frac12}
		=
		\frac12.
		\tag{1}
		\]
		
		Необходимо найти
		\[
		(\omega_1,\omega_2,t_1,t_2,\beta).
		\]
		
		Будем считать $x$ заданным параметром.
		
		\subsection*{1. Дополнительное уравнение для $t_1,t_2$}
		
		Из (1):
		\[
		\frac{1-t_2}{t_2-t_1}=\frac12.
		\]
		
		Следовательно,
		\[
		2(1-t_2)=t_2-t_1.
		\]
		
		Отсюда
		\[
		2-2t_2=t_2-t_1,
		\]
		поэтому
		\[
		\boxed{t_1=3t_2-2},
		\tag{2}
		\]
		или
		\[
		\boxed{t_2=\frac{t_1+2}{3}}.
		\tag{3}
		\]
		
		Кроме того, поскольку исходное условие содержит знаменатель
		$t_2-t_1$, необходимо
		\[
		\boxed{t_2\neq t_1}.
		\tag{4}
		\]
		
		\subsection*{2. Разбиение векторного уравнения на скалярные}
		
		Получаем четыре уравнения:
		\begin{align}
			0&=\omega_1-\omega_2+(1-\omega_1-\omega_2),
			\tag{5}\\
			\beta&=\omega_1t_1-\omega_2t_2+(1-\omega_1-\omega_2),
			\tag{6}\\
			2\beta x&=\omega_1t_1^2-\omega_2t_2^2+(1-\omega_1-\omega_2),
			\tag{7}\\
			3\beta x^2&=\omega_1t_1^3-\omega_2t_2^3+(1-\omega_1-\omega_2).
			\tag{8}
		\end{align}
		
		\subsection*{3. Нахождение $\omega_2$}
		
		Из (5):
		\[
		0=\omega_1-\omega_2+1-\omega_1-\omega_2.
		\]
		
		Следовательно,
		\[
		1-2\omega_2=0,
		\]
		откуда
		\[
		\boxed{\omega_2=\frac12}.
		\tag{9}
		\]
		
		Поэтому
		\[
		1-\omega_1-\omega_2
		=
		\frac12-\omega_1.
		\tag{10}
		\]
		
		\subsection*{4. Нахождение $\beta$}
		
		Подставляем (9), (10) в (6):
		\[
		\beta
		=
		\omega_1t_1
		-\frac{t_2}{2}
		+\frac12-\omega_1.
		\]
		
		То есть
		\[
		\beta
		=
		\omega_1(t_1-1)+\frac{1-t_2}{2}.
		\tag{11}
		\]
		
		Используем
		\[
		t_2=\frac{t_1+2}{3}.
		\]
		
		Тогда
		\[
		\frac{1-t_2}{2}
		=
		\frac12
		\left(
		1-\frac{t_1+2}{3}
		\right)
		=
		\frac{1-t_1}{6}.
		\]
		
		Следовательно,
		\[
		\beta
		=
		\omega_1(t_1-1)+\frac{1-t_1}{6}.
		\]
		
		Получаем
		\[
		\boxed{
			\beta=(t_1-1)
			\left(
			\omega_1-\frac16
			\right)
		}.
		\tag{12}
		\]
		
		\subsection*{5. Используем третье скалярное уравнение}
		
		Рассмотрим (7):
		\[
		2\beta x
		=
		\omega_1t_1^2
		-\frac12t_2^2
		+\frac12-\omega_1.
		\]
		
		То есть
		\[
		2\beta x
		=
		\omega_1(t_1^2-1)
		+\frac{1-t_2^2}{2}.
		\tag{13}
		\]
		
		Поскольку
		\[
		t_2=\frac{t_1+2}{3},
		\]
		имеем
		\[
		t_2^2
		=
		\frac{(t_1+2)^2}{9}.
		\]
		
		Поэтому
		\[
		\frac{1-t_2^2}{2}
		=
		\frac12
		\left(
		1-\frac{(t_1+2)^2}{9}
		\right).
		\]
		
		Раскрываем:
		\[
		\frac{1-t_2^2}{2}
		=
		\frac{9-(t_1^2+4t_1+4)}{18},
		\]
		то есть
		\[
		\frac{1-t_2^2}{2}
		=
		-\frac{t_1^2+4t_1-5}{18}.
		\]
		
		Разложим числитель:
		\[
		t_1^2+4t_1-5=(t_1-1)(t_1+5).
		\]
		
		Поэтому
		\[
		2\beta x
		=
		\omega_1(t_1-1)(t_1+1)
		-
		\frac{(t_1-1)(t_1+5)}{18}.
		\]
		
		Следовательно,
		\[
		2\beta x
		=
		(t_1-1)
		\left[
		\omega_1(t_1+1)
		-\frac{t_1+5}{18}
		\right].
		\tag{14}
		\]
		
		С другой стороны, из (12):
		\[
		2\beta x
		=
		2x(t_1-1)
		\left(
		\omega_1-\frac16
		\right).
		\tag{15}
		\]
		
		Поскольку $t_1\neq1$ в допустимом решении, сокращаем $t_1-1$:
		\[
		2x
		\left(
		\omega_1-\frac16
		\right)
		=
		\omega_1(t_1+1)
		-\frac{t_1+5}{18}.
		\]
		
		Умножаем на $18$:
		\[
		36x\omega_1-6x
		=
		18\omega_1(t_1+1)-t_1-5.
		\]
		
		Отсюда
		\[
		18\omega_1(t_1+1-2x)
		=
		t_1+5-6x.
		\]
		
		Следовательно,
		\[
		\boxed{
			\omega_1=
			\frac{t_1+5-6x}
			{18(t_1+1-2x)}
		}.
		\tag{16}
		\]
		
		\subsection*{6. Используем четвёртое уравнение}
		
		Теперь используем (8):
		\[
		3\beta x^2
		=
		\omega_1t_1^3
		-\frac12t_2^3
		+\frac12-\omega_1.
		\]
		
		Перепишем:
		\[
		3\beta x^2
		=
		\omega_1(t_1^3-1)
		+\frac{1-t_2^3}{2}.
		\tag{17}
		\]
		
		Так как
		\[
		t_2=\frac{t_1+2}{3},
		\]
		то
		\[
		t_2^3
		=
		\frac{(t_1+2)^3}{27}.
		\]
		
		Следовательно,
		\[
		\frac{1-t_2^3}{2}
		=
		\frac{27-(t_1+2)^3}{54}.
		\]
		
		Раскрывая куб:
		\[
		(t_1+2)^3
		=
		t_1^3+6t_1^2+12t_1+8,
		\]
		получаем
		\[
		\frac{1-t_2^3}{2}
		=
		-\frac{
			t_1^3+6t_1^2+12t_1-19
		}{54}.
		\]
		
		Также
		\[
		t_1^3-1
		=
		(t_1-1)(t_1^2+t_1+1).
		\]
		
		Подставляем (12), (16) в (17). После раскрытия скобок и
		приведения подобных членов получается
		\[
		\boxed{
			(t_1-1)^2
			\left(
			t_1^2+(6-8x)t_1+9x^2-10x+2
			\right)
			=0.
		}
		\tag{18}
		\]
		
		Но $t_1=1$ недопустимо: из (3) тогда
		$t_2=1$, а исходная дробь
		\[
		\frac{1-t_2}{t_2-t_1}
		\]
		имеет нулевой знаменатель.
		
		Поэтому
		\[
		\boxed{
			t_1^2+(6-8x)t_1+9x^2-10x+2=0.
		}
		\tag{19}
		\]
		
		\subsection*{7. Нахождение $t_1$}
		
		Дискриминант:
		\[
		D
		=
		(6-8x)^2
		-
		4(9x^2-10x+2).
		\]
		
		Раскрываем:
		\[
		D
		=
		64x^2-96x+36
		-36x^2+40x-8.
		\]
		
		Следовательно,
		\[
		D=28x^2-56x+28,
		\]
		то есть
		\[
		\boxed{
			D=28(x-1)^2
		}.
		\tag{20}
		\]
		
		Поэтому
		\[
		t_1
		=
		\frac{8x-6\pm\sqrt{28(x-1)^2}}{2}.
		\]
		
		Введём
		\[
		s\in\{-1,1\}.
		\]
		
		Тогда
		\[
		\boxed{
			t_1
			=
			4x-3+s\sqrt{7}\,(1-x)
		}.
		\tag{21}
		\]
		
		\subsection*{8. Нахождение $t_2$}
		
		Используем (3):
		\[
		t_2=\frac{t_1+2}{3}.
		\]
		
		Подставляем (21):
		\[
		t_2
		=
		\frac{
			4x-3+s\sqrt{7}(1-x)+2
		}{3}.
		\]
		
		Следовательно,
		\[
		\boxed{
			t_2
			=
			\frac{
				4x-1+s\sqrt{7}\,(1-x)
			}{3}
		}.
		\tag{22}
		\]
		
		\subsection*{9. Нахождение $\omega_1$}
		
		Используем (16):
		\[
		\omega_1
		=
		\frac{t_1+5-6x}
		{18(t_1+1-2x)}.
		\]
		
		Подставляем
		\[
		t_1=4x-3+s\sqrt{7}(1-x).
		\]
		
		В числителе:
		\[
		t_1+5-6x
		=
		1-2x+s\sqrt{7}(1-x).
		\]
		
		В знаменателе:
		\[
		t_1+1-2x
		=
		2x-2+s\sqrt{7}(1-x).
		\]
		
		После упрощения получаем
		\[
		\boxed{
			\omega_1
			=
			\frac{11+4s\sqrt{7}}{54}
		}.
		\tag{23}
		\]
		
		При этом
		\[
		\boxed{
			\omega_2=\frac12
		}.
		\tag{24}
		\]
		
		\subsection*{10. Нахождение $\beta$}
		
		Из (12):
		\[
		\beta
		=
		(t_1-1)
		\left(
		\omega_1-\frac16
		\right).
		\]
		
		Имеем
		\[
		t_1-1
		=
		4x-4+s\sqrt{7}(1-x)
		=
		(x-1)(4-s\sqrt{7}),
		\]
		а
		\[
		\omega_1-\frac16
		=
		\frac{11+4s\sqrt{7}-9}{54}
		=
		\frac{2+4s\sqrt{7}}{54}.
		\]
		
		После упрощения:
		\[
		\boxed{
			\beta
			=
			\frac{(1-x)(10-7s\sqrt{7})}{27}
		}.
		\tag{25}
		\]
		
		\subsection*{11. Окончательный ответ}
		
		Таким образом, при
		\[
		s=\pm1
		\]
		получаем две ветви решения:
		\[
		\boxed{
			\omega_1=
			\frac{11+4s\sqrt{7}}{54}
		},
		\]
		\[
		\boxed{
			\omega_2=\frac12
		},
		\]
		\[
		\boxed{
			t_1=
			4x-3+s\sqrt{7}(1-x)
		},
		\]
		\[
		\boxed{
			t_2=
			\frac{4x-1+s\sqrt{7}(1-x)}{3}
		},
		\]
		\[
		\boxed{
			\beta=
			\frac{(1-x)(10-7s\sqrt{7})}{27}
		}.
		\]
		
		Или компактно:
		\[
		\boxed{
			(\omega_1,\omega_2,t_1,t_2,\beta)
			=
			\left(
			\frac{11+4s\sqrt{7}}{54},
			\frac12,
			4x-3+s\sqrt{7}(1-x),
			\frac{4x-1+s\sqrt{7}(1-x)}{3},
			\frac{(1-x)(10-7s\sqrt{7})}{27}
			\right),
			\qquad s=\pm1.
		}
		\tag{26}
		\]
		
		При этом необходимо
		\[
		\boxed{x\neq1}.
		\]
		
		Действительно, при $x=1$ обе ветви дают
		\[
		t_1=t_2=1,
		\]
		а исходное условие
		\[
		\frac{1-t_2}{t_2-t_1}=\frac12
		\]
		становится недопустимым из-за нулевого знаменателя.
		
		Итак, для $x\neq1$ существуют две ветви решения, соответствующие
		$s=\pm1$.
		
	\newpage
	\section{Нечебышёвский, внутренности}
	$$
	\beta
	\begin{pmatrix}
		0\\
		1\\
		2x\\
		3x^2
	\end{pmatrix}
	=
	\omega_1
	\begin{pmatrix}
		1\\
		-1\\
		1\\
		-1
	\end{pmatrix}
	-\omega_2
	\begin{pmatrix}
		1\\
		t_2\\
		t_2^2\\
		t_2^3
	\end{pmatrix}
	+
	(1-\omega_1-\omega_2)
	\begin{pmatrix}
		1\\
		1\\
		1\\
		1
	\end{pmatrix}.
	$$
	
	Считаем \(x\) заданным параметром. Приравнивая координаты, получаем систему
	
	$$
	\begin{cases}
		0=\omega_1-\omega_2+1-\omega_1-\omega_2,\\
		\beta=-\omega_1-\omega_2t_2+1-\omega_1-\omega_2,\\
		2x\beta=\omega_1-\omega_2t_2^2+1-\omega_1-\omega_2,\\
		3x^2\beta=-\omega_1-\omega_2t_2^3+1-\omega_1-\omega_2.
	\end{cases}
	$$
	
	Из первого уравнения сразу получаем
	
	$$
	0=1-2\omega_2,
	$$
	
	откуда
	
	$$
	\boxed{\omega_2=\frac12}.
	$$
	
	Второе уравнение принимает вид
	
	$$
	\beta
	=
	1-2\omega_1-\frac{t_2+1}{2},
	$$
	
	следовательно,
	
	$$
	\boxed{
		\omega_1=\frac{1-t_2-2\beta}{4}
	}.
	$$
	
	Рассмотрим третье уравнение:
	
	$$
	2x\beta
	=
	\omega_1-\frac{t_2^2}{2}
	+1-\omega_1-\frac12.
	$$
	
	Таким образом,
	
	$$
	2x\beta=\frac{1-t_2^2}{2},
	$$
	
	и, при \(x\neq0\),
	
	$$
	\boxed{
		\beta=\frac{1-t_2^2}{4x}
	}.
	$$
	
	Теперь рассмотрим четвёртое уравнение:
	
	$$
	3x^2\beta
	=
	-\omega_1-\frac{t_2^3}{2}
	+1-\omega_1-\frac12.
	$$
	
	Подставляя выражение для \(\omega_1\), получаем
	
	$$
	3x^2\beta
	=
	\frac{1-t_2^3}{2}
	-\frac{1-t_2-2\beta}{2}.
	$$
	
	После упрощения:
	
	$$
	3x^2\beta
	=
	\frac{t_2-t_2^3+2\beta}{2}.
	$$
	
	Следовательно,
	
	$$
	3x^2\beta
	=
	\frac{t_2(1-t_2^2)}{2}+\beta,
	$$
	
	то есть
	
	$$
	\beta(3x^2-1)
	=
	\frac{t_2(1-t_2^2)}{2}.
	$$
	
	Используя
	
	$$
	\beta=\frac{1-t_2^2}{4x},
	$$
	
	получаем
	
	$$
	\frac{1-t_2^2}{4x}(3x^2-1)
	=
	\frac{t_2(1-t_2^2)}{2}.
	$$
	
	Отсюда
	
	$$
	(1-t_2^2)
	\left(
	\frac{3x^2-1}{4x}
	-\frac{t_2}{2}
	\right)=0.
	$$
	
	Таким образом, получаем две ветви решений.
	
	\subsection*{1. Случай \(t_2=\pm1\)}
	
	Если
	
	$$
	1-t_2^2=0,
	$$
	
	то
	
	$$
	t_2=\pm1.
	$$
	
	При этом
	
	$$
	\beta=0.
	$$
	
	Из
	
	$$
	\omega_1=\frac{1-t_2-2\beta}{4}
	$$
	
	получаем:
	
	$$
	\boxed{
		t_2=1,\qquad
		\beta=0,\qquad
		\omega_1=0,\qquad
		\omega_2=\frac12
	}
	$$
	
	и
	
	$$
	\boxed{
		t_2=-1,\qquad
		\beta=0,\qquad
		\omega_1=\frac12,\qquad
		\omega_2=\frac12
	}.
	$$
	
	\subsection*{2. Случай \(t_2\neq\pm1\)}
	
	Тогда можно сократить множитель \(1-t_2^2\):
	
	$$
	\frac{3x^2-1}{4x}
	=
	\frac{t_2}{2}.
	$$
	
	Отсюда
	
	$$
	\boxed{
		t_2=\frac{3x^2-1}{2x}
	}.
	$$
	
	Подставляем это выражение в
	
	$$
	\beta=\frac{1-t_2^2}{4x}.
	$$
	
	Получаем
	
	$$
	\beta
	=
	\frac{
		1-\left(\frac{3x^2-1}{2x}\right)^2
	}{4x}.
	$$
	
	После упрощения:
	
	$$
	\boxed{
		\beta
		=
		-\frac{(x^2-1)(9x^2-1)}{16x^3}
	}.
	$$
	
	Для \(\omega_1\):
	
	$$
	\omega_1
	=
	\frac{1-t_2-2\beta}{4}.
	$$
	
	После подстановки и упрощения:
	
	$$
	\boxed{
		\omega_1
		=
		-\frac{(x-1)^3(3x+1)}{32x^3}
	}.
	$$
	
	При этом
	
	$$
	\boxed{\omega_2=\frac12}.
	$$
	
	Таким образом, третья ветвь имеет вид
	
	$$
	\boxed{
		\begin{aligned}
			t_2&=\frac{3x^2-1}{2x},\\
			\beta&=-\frac{(x^2-1)(9x^2-1)}{16x^3},\\
			\omega_1&=-\frac{(x-1)^3(3x+1)}{32x^3},\\
			\omega_2&=\frac12.
		\end{aligned}
	}
	$$
	
	Итак, при \(x\neq0\) все решения системы задаются тремя ветвями:
	
	$$
	\boxed{
		\begin{aligned}
			&t_2=1,\qquad
			\beta=0,\qquad
			\omega_1=0,\qquad
			\omega_2=\frac12;
			\\[2mm]
			&t_2=-1,\qquad
			\beta=0,\qquad
			\omega_1=\frac12,\qquad
			\omega_2=\frac12;
			\\[2mm]
			&t_2=\frac{3x^2-1}{2x},\qquad
			\beta=-\frac{(x^2-1)(9x^2-1)}{16x^3},\qquad
			\omega_1=-\frac{(x-1)^3(3x+1)}{32x^3},\qquad
			\omega_2=\frac12.
		\end{aligned}
	}
	$$
	
	При \(x=\pm1,\pm\frac13\) третья ветвь совпадает с одной из первых двух.
	
	
	
\end{document}

